\section{The adjoint-defect receiver and its actual invariant-cycle
cross}\label{the-adjoint-defect-receiver-and-its-actual-invariant-cycle-cross}

Complete local derivation. Proof locators \textbf{ADC0--ADC9}. This note
calculates the invariant-cycle cross of the existing receiver of the
positive-adjoint discrepancy. Its attaching map is retained; neither a
new coefficient system nor a purity hypothesis is inserted.

\subsection{ADC0. Construction stage, correction audit, and
inputs}\label{adc0.-construction-stage-correction-audit-and-inputs}

The supporting datum is \(\tau\langle Z_1;\mathrm{no}\ Z_2\rangle\). All
operations here take place on the coefficient spaces after
complete-history arithmetic reconstruction. No addition, coordinate,
metric, midpoint, or extra parity is assigned to the support. The two
arithmetic branches retain their own recovered counters.

Before this calculation the connected
\nolinkurl{USER_ARGUMENT_RECONSTRUCTION.md,} its correction-precedence
table, \nolinkurl{READ_FIRST_USER_CONSTRUCTION.md,} the verbatim WU062
passage, and the WU064--WU065 amendment were read. In particular, an
obstruction in a specified coefficient receiver must be connected by
actual maps to the other programme objects; an independently selected
extension cannot replace the desired lifting problem. WU065 also
requires attempting the next calculation using all established results
when a preceding result does not give the sought conclusion.

The actual inputs are RGR0--RGR12 in
\nolinkurl{CC_RESIDUE_GYSIN_RECEIVER_INDEPENDENT.md;} LNC0--LNC11 in
\nolinkurl{CC_LOCAL_NEARBY_CYCLES_AND_SPECIALIZATION.md;} LVD0--LVD9 in
\nolinkurl{CC_LOCAL_VANISHING_DUAL_ARROWS.md;} and DC0--DC12 in
\nolinkurl{DELIGNE_INVARIANT_CYCLE_QUOTIENT.md.} RGR constructs the
source map and current cone. LNC constructs nearby cycles, the exact
local cross, and the proper-current/cover signs. DC proves the two
weight directions and the exact quotient for Deligne's distinct
geometric objects. The positive-adjoint operator is the actual
GTAH0--GTAH1 operator in
\nolinkurl{GEOMETRIC_TRANSFER_POSITIVE_ADJOINT_DEFECT.md.} No source
theorem for finite-rank constructible sheaves is applied to infinite
coefficients without the explicit local calculation.

Sheaf derived statements use sheaves of complex vector spaces, and all
specified coefficient/current maps are continuous. Kernels and quotients
retain their actual topologies; no closed image is inferred from an
algebraic quotient. Continuous duals are distinguished from unrestricted
algebraic duals.

\subsection{ADC1. The full source and the particular attaching
map}\label{adc1.-the-full-source-and-the-particular-attaching-map}

Use the original source spaces and quotient \[
S=\{f\in\mathcal S(\mathbb R;\mathbb C):f(-v)=f(v),\
 f(0)=0,\ \int_{\mathbb R}f(v)\,dv=0\},
\] \[
A=\{a\in C^\infty(\mathbb R_{>0}):
\sup_{u>0}(u^N+u^{-N})|(u\partial_u)^ja(u)|<\infty\
\text{for every }N,j\},\quad
\Sigma f(u)=2\sum_{n\geq1}f(nu),
\] \[
J=\Sigma S,\qquad Q=A/J,\qquad q:A\to Q,\qquad
B=A'_\beta,\quad j=q':Q'_\beta\hookrightarrow B .
\tag{ADC1.1}
\] SSI/ESI give closed \(J\) and the continuous inverse of
\(\Sigma:S\to J\). SDT/CSD give the strong embedding \(j(Q')=J^\perp\).
The prime denotes the continuous complex-linear dual, and \(\beta\)
denotes uniform convergence on bounded subsets.

The complete original transform is \[
\Theta a(s)=\frac12\int_0^\infty a(u)u^s\frac{du}{u},\qquad
\Theta^{-1}F(u)=\frac{u^{-1/2}}{\pi}
\int_{\mathbb R}F(1/2+it)u^{-it}\,dt.
\tag{ADC1.2}
\] It identifies \(J\) with the entire strip-Schwartz ideal of full
vanishing jets at the original nontrivial zeros, with orders \(m_\rho\).
The retained divisor is \[
f_0(v)=\frac{\pi}{2}v^2(2\pi v^2-3)e^{-\pi v^2},\qquad
F_0(s)=\Theta\Sigma f_0(s)
=\frac{s(s-1)}8\pi^{-s/2}\Gamma(s/2)\zeta(s).
\tag{ADC1.3}
\] Its exceptional values remain \[
F_0(0)=F_0(1)=\frac18,\quad F_0(-1)=F_0(2)=\frac{\pi}{24},
\quad
F_0(-2k)=\frac{k(2k+1)(-1)^k\pi^k}{2\,k!}\zeta'(-2k)\ne0.
\tag{ADC1.4}
\] At \(\rho+t\), the full local factor is \[
F_0(\rho+t)=t^{m_\rho}
\frac{(\rho+t)(\rho+t-1)}8\pi^{-(\rho+t)/2}
\Gamma((\rho+t)/2)\frac{\zeta(\rho+t)}{t^{m_\rho}}.
\tag{ADC1.5}
\] Every derivative of this factor is retained. The original arithmetic
is \[
\zeta(s)=1+\sum_{n\ge2}n^{-s}=\prod_p(1-p^{-s})^{-1},\quad
-\frac{\zeta'(s)}{\zeta(s)}
=\sum_p\sum_{k\ge1}(\log p)p^{-ks}\quad(\Re s>1).
\tag{ADC1.6}
\] Thus the unit and all prime-power repetitions remain in the source.
No modified zero divisor is substituted.

Let \(\mathscr Z\) be the distinct original nontrivial zeros,
\(\rho^\#=1-\overline\rho\), and \[
H=\ell^2(\mathscr Z,m),\quad
\langle x,y\rangle_+=\sum_\rho m_\rho x_\rho\overline{y_\rho},
\quad(\mathsf Jy)_\rho=y_{\rho^\#}.
\] The actual continuous map \(E:Q\to H\) is
\(E(F)_\rho=(\Theta_QF)(\rho)\). Its image contains every finite value
vector and is dense. Its full higher-jet kernel is retained in \(Q\).
Put \(C=\overline H\), with its conjugate complex structure. For a fixed
recovered parameter \(r>1\), \[
(T_ry)_\rho=r^\rho y_\rho,\quad U_r=rT_{1/r},\quad
D_r=T_r^*-U_r,\qquad
d_r(\rho):=e^{-i\Im\rho\log r}
 \bigl(r^{\Re\rho}-r^{1-\Re\rho}\bigr).
\tag{ADC1.7}
\] Then \((D_ry)_\rho=d_r(\rho)y_\rho\). These are bounded operators by
the known open strip. An integer \(r=n\) is a geometric cover degree; a
noninteger \(r\) is only a coefficient parameter.

The precise attaching map is \[
A_H(y)(F)=\langle EF,\mathsf Jy\rangle_+,\quad
\sigma_r:C\to Q'_\beta,\quad
\sigma_r(\overline y)=A_H(D_ry),\qquad
a_r=j\sigma_r:C\to B.
\tag{ADC1.8}
\] \(A_H\) is anti-linear, so \(\sigma_r\) is complex-linear. For a
bounded test set \(K\subset Q\), \[
\sup_{F\in K}|\sigma_r(\overline y)(F)|
\leq \sup_{F\in K}\|EF\|_+\,\|D_r\|\,\|y\|_+,
\tag{ADC1.9}
\] proving strong continuity and absolute convergence. Density of
\(E(Q)\) proves injectivity of \(A_H\); \(j\) is injective. Consequently
\[
\ker a_r=\ker\sigma_r=\overline{\ker D_r}
 =:C_{\rm crit},
\tag{ADC1.10}
\] where \(C_{\rm crit}\) consists exactly of coordinates supported at
the actual critical-line zeros. The equality follows from strict
monotonicity of \(r^x\) in (ADC1.7); it asserts neither existence nor
nonexistence of an off-line zero.

For comparison with the original residue realization, if
\(\zeta(\rho+t)=t^{m_\rho}u_\rho(t)\), the functional in (ADC1.8) is the
strong limit of the original-\(\zeta\) residues with finite numerators
\[
h_{B_0}(y)=\sum_{\rho\in B_0}
m_\rho(-1)^{m_\rho-1}u_\rho(0)\,
\overline{(D_ry)_{\rho^\#}}\,
e_{1-\rho,m_\rho-1},
\quad
\mathcal R(F,h)=\sum_\rho\operatorname{Res}_{s=\rho}
\frac{F(s)h(1-s)}{\zeta(s)}\,ds.
\tag{ADC1.11}
\] The \(e\)'s are RTT's full-jet isolators. Its residue calculation
gives the exact coefficient of the value functional; (ADC1.9) controls
the strong limit. Neither multiplicity nor the leading original-zeta
unit is omitted.

\subsection{ADC2. The current cone, endpoints, and its four local
functors}\label{adc2.-the-current-cone-endpoints-and-its-four-local-functors}

Fix a pole disk with positive coordinate \(z\) at zero or \(1/z\) at
infinity. Let \(\mathscr T_A\) be the continuous compact-test current
complex, \[
\mathscr T_A^n(U)=(\mathcal E_c^{-n}(U;A))',
\quad dT(\phi)=(-1)^{n+1}T(d\phi),\quad -2\leq n\leq0.
\] Its constants represent \(\underline B[2]\). Positive Dirac and
angular currents satisfy \[
\delta_p\lambda(f)=\lambda(f(p)),\quad
h_p\lambda(\omega)=\lambda\left(\int\frac{d\arg z_p}{2\pi}\wedge\omega\right),
\quad dh_p\lambda=\delta_p\lambda.
\tag{ADC2.1}
\] The last equality is Stokes with positive small-circle integral one.
All \(A\)-valued integrals are formed before the functional is applied.

The actual receiver is \[
\mathscr E_r=\operatorname{Cone}(-\delta_pa_r)[-1].
\] In degrees \(-1,0,1\) its terms and differentials are \[
\mathscr T_A^{-2},\quad\mathscr T_A^{-1}\oplus i_*C,\quad
\mathscr T_A^0,\qquad
d^{-1}b=(-d_Tb,0),\quad
d^0(b,x)=-d_Tb+\delta_pa_rx.
\tag{ADC2.2}
\] It is RGR's actual homotopy pullback; the triangle is \[
\underline B[1]\longrightarrow\mathscr E_r
\longrightarrow i_*C\xrightarrow{\mathrm{Gys}_B\,a_r}\underline B[2].
\tag{ADC2.3}
\] On a disk the constant and angular representatives give the
continuous cohomology models \[
i^*\mathscr E_r\simeq B[1]\oplus C,\qquad
R\Gamma(D^*,\mathscr E_r)\simeq B[1]\oplus B[\vartheta_p],
\] \[
\lambda\longmapsto c_\lambda\quad(\deg-1),\qquad
x\longmapsto(h_pa_rx,x)\quad(\deg0),
\quad \mathrm{res}=(1_B,a_r).
\tag{ADC2.4}
\] Closedness of the latter representative follows from (ADC2.1); the
compact-test contraction removes the remaining current terms. These are
continuous comparisons on fixed disks, and their cohomology
identifications agree as disks shrink, as in CSD/RGR.

Pulling the punctured constant system to the universal cover and
contracting leaves \[
R\psi\mathscr E_r=B[1],\quad T=1,\qquad
\operatorname{sp}:B[1]\oplus C\longrightarrow B[1]
\text{ is projection}.
\tag{ADC2.5}
\] Its vanishing cone is \(R\phi\mathscr E_r=C[1]\). Equivalently the
shifted functors \(\Psi=R\psi[-1]\), \(\Phi=R\phi[-1]\) have \[
\Psi\mathscr E_r=B,\quad \Phi\mathscr E_r=C,\qquad
\mathrm{can}=0,\quad\mathrm{var}=a_r.
\tag{ADC2.6}
\] The positive angular convention fixes the plus in variation; it is
the second restriction component in (ADC2.4), not a sign deduced from
\(1-T=0\).

The restriction fibre first has coefficient complex \[
B\xrightarrow{b\mapsto(0,b)}C\oplus B
\xrightarrow{(x,b)\mapsto a_rx}B
\quad\text{in degrees }-1,0,1.
\tag{ADC2.7}
\] Cancel its identity pair to obtain \[
i^!\mathscr E_r\simeq[C\xrightarrow{+a_r}B]
\quad(\deg0,1).
\tag{ADC2.8}
\] The target is positively oriented: in the shifted current
differential, \(d h_pb=-\delta_pb\), so the fibre differential gives
\(d(h_pb,0)=(-\delta_pb,h_pb)\). The positive Dirac class is therefore
the positive angular boundary coordinate. This explains why (ADC2.8) has
\(+a_r\), whereas a literal standard shifted cone on \(a_r\) has
differential \(-a_r\); negating its target gives the displayed
comparison.

Both source endpoint lines remain as in RGR10: \[
\mathscr E_{r,\mathrm{full}}=\mathscr E_r\oplus i_*E_p'[-1].
\tag{ADC2.9}
\] Their stalk and costalk are \(E_p'[-1]\), their nearby complex is
zero, and their unshifted vanishing complex is \(E_p'[0]\). Thus they
appear in special and support degree \(1\), and vanishing degree \(0\),
separately from (ADC2.5)--(ADC2.8).

Here \(E_+=\mathbb C c_0\oplus\mathbb C c_1\) and
\(E_-=\mathbb C d_0\oplus\mathbb C d_1\) retain respectively the
original value-at-zero and integral channels in the two source charts.
Their original coefficient actions are \[
\rho_+(t)|_{E_+}=\operatorname{diag}(1,t),\qquad
\rho_-(t)|_{E_-}=\operatorname{diag}(t,1).
\tag{ADC2.10}
\] These are the endpoint components of the complete source maps in
(LNC5.2); transposition keeps each indicated coordinate label.

\subsection{ADC3. The exact invariant-cycle cross of this
receiver}\label{adc3.-the-exact-invariant-cycle-cross-of-this-receiver}

Write \[
S_r=a_r(C)\subset B,\qquad N_r=B/S_r.
\tag{ADC3.1}
\] These are the actual image and algebraic quotient with their stated
topologies. No closedness is asserted. From (ADC2.8), the full support
groups are \[
H^0_p(\mathscr E_{r,\mathrm{full}})=C_{\rm crit},\qquad
H^1_p(\mathscr E_{r,\mathrm{full}})=N_r\oplus E_p'.
\tag{ADC3.2}
\] Let \(V_i=H^i(R\psi\mathscr E_{r,\mathrm{full}})\),
\(A_i^{\rm sp}=H^i(i^*\mathscr E_{r,\mathrm{full}})\),
\(B_i^{\rm punct}=H^i(D^*,\mathscr E_{r,\mathrm{full}})\). Define the
actual invariants \(C_i^{\rm inv}=V_i^T\), angular coinvariants
\(K_i=V_{i-1,T}[\vartheta_p]\), and \(O_i=H^{i+1}_p\). The only nearby
group is \(V_{-1}=B\), with monodromy identity. Thus the entire nonzero
cross is:

{\def\LTcaptype{none} % do not increment counter
\begin{longtable}[]{@{}
  >{\raggedright\arraybackslash}p{(\linewidth - 14\tabcolsep) * \real{0.1250}}
  >{\raggedright\arraybackslash}p{(\linewidth - 14\tabcolsep) * \real{0.1250}}
  >{\raggedright\arraybackslash}p{(\linewidth - 14\tabcolsep) * \real{0.1250}}
  >{\raggedright\arraybackslash}p{(\linewidth - 14\tabcolsep) * \real{0.1250}}
  >{\raggedright\arraybackslash}p{(\linewidth - 14\tabcolsep) * \real{0.1250}}
  >{\raggedright\arraybackslash}p{(\linewidth - 14\tabcolsep) * \real{0.1250}}
  >{\raggedright\arraybackslash}p{(\linewidth - 14\tabcolsep) * \real{0.1250}}
  >{\raggedright\arraybackslash}p{(\linewidth - 14\tabcolsep) * \real{0.1250}}@{}}
\toprule\noalign{}
\begin{minipage}[b]{\linewidth}\raggedright
\(i\)
\end{minipage} & \begin{minipage}[b]{\linewidth}\raggedright
\(A_i^{\rm sp}\)
\end{minipage} & \begin{minipage}[b]{\linewidth}\raggedright
\(K_i\)
\end{minipage} & \begin{minipage}[b]{\linewidth}\raggedright
\(B_i^{\rm punct}\)
\end{minipage} & \begin{minipage}[b]{\linewidth}\raggedright
\(C_i^{\rm inv}\)
\end{minipage} & \begin{minipage}[b]{\linewidth}\raggedright
\(O_i\)
\end{minipage} & \begin{minipage}[b]{\linewidth}\raggedright
\(\alpha_i\)
\end{minipage} & \begin{minipage}[b]{\linewidth}\raggedright
\(\partial_i\)
\end{minipage} \\
\midrule\noalign{}
\endhead
\bottomrule\noalign{}
\endlastfoot
\(-1\) & \(B\) & \(0\) & \(B\) & \(B\) & \(C_{\rm crit}\) & \(1\) &
\(0\) \\
\(0\) & \(C\) & \(B[\vartheta_p]\) & \(B[\vartheta_p]\) & \(0\) &
\(N_r\oplus E_p'\) & \(a_r\) & \(b\mapsto([b],0)\) \\
\(1\) & \(E_p'\) & \(0\) & \(0\) & \(0\) & \(0\) & \(0\) & \(0\) \\
\end{longtable}
}

The row
\(0\to K_i\xrightarrow{j_i}B_i^{\rm punct}\to C_i^{\rm inv}\to0\) has
identity \(\pi_{-1}\) and identity \(j_0\). The complete remaining
localization part is \[
0\to C_{\rm crit}\to C\xrightarrow{a_r}B
\xrightarrow{b\mapsto([b],0)}N_r\oplus E_p'
\xrightarrow{\mathrm{pr}_{E_p'}}E_p'\to0.
\tag{ADC3.3}
\] Every sign follows from (ADC2.7)--(ADC2.8).

It follows by these actual maps, without a weight argument, that \[
\boxed{\operatorname{im}\partial_i/\partial_i j_iK_i=0
\quad\text{for every }i.}
\tag{ADC3.4}
\] At degree \(-1\) the boundary is zero. At degree zero the coinvariant
injection is identity, so its boundary image is all \(N_r\oplus0\). The
quotient of the whole receiver is still \[
O_0/\partial_0j_0K_0=E_p',
\tag{ADC3.5}
\] not zero when endpoints are retained. This is exactly the remaining
specialization kernel in degree one. Equations (ADC3.4)--(ADC3.5)
distinguish the obstruction image from the whole quotient receiver.

The vanishing of (ADC3.4) does not imply the attaching map vanished: it
holds for the present \(a_r\) before any conclusion about \(D_r\). The
stronger conclusion of DC8 has a different test here: \[
\ker(\partial_0j_0)=S_r,\qquad
K_0\xrightarrow{\partial_0j_0}\operatorname{im}\partial_0=N_r
\text{ is injective }\Longleftrightarrow a_r=0
\Longleftrightarrow D_r=0.
\tag{ADC3.6}
\] The actual source subquotient in DC5.5 is \[
0\to C_{\rm crit}\to C\xrightarrow{a_r}S_r\to0.
\tag{ADC3.7}
\] The map \(C/C_{\rm crit}\to S_r\) is a continuous algebraic
bijection. We do not assume its inverse continuous for the ambient
subspace topology of \(S_r\).

There is also an exact calculation of the stronger DC9 candidate: \[
\frac{A_0^{\rm sp}}{\operatorname{im}(H^0_p\to A_0^{\rm sp})}
=C/C_{\rm crit}\longrightarrow C_0^{\rm inv}=0.
\tag{ADC3.8}
\] Its kernel is the entire displayed source quotient, identified
algebraically with \(S_r\). Thus the actual next object that must be
controlled after the automatically zero invariant-cycle obstruction is
\(S_r\), not the already zero quotient (ADC3.4).

\subsection{ADC4. The exact relation to the original coefficient
quotient}\label{adc4.-the-exact-relation-to-the-original-coefficient-quotient}

Write \(S_r^Q=\sigma_r(C)\subset Q'\). Since \(j\) is injective,
\(S_r=jS_r^Q\). The full continuous-transpose coefficient row is \[
0\to Q'\xrightarrow jB\xrightarrow{k'}J'\to0,
\tag{ADC4.1}
\] where \(k:J\hookrightarrow A\). Exactness is the quotient definition
and Hahn--Banach on closed \(J\); no strong-open mapping theorem is
used. It induces the exact row \[
0\to Q'/S_r^Q\xrightarrow{[\lambda]\mapsto[j\lambda]}
N_r\xrightarrow{[b]\mapsto k'b}J'\to0.
\tag{ADC4.2}
\] For injectivity, \(j\lambda\in S_r=jS_r^Q\) implies
\(\lambda\in S_r^Q\). The middle kernel consists of classes with
\(k'b=0\), hence \(b=j\lambda\). Surjectivity is that of \(k'\). These
prove every exactness assertion.

The row is exactly the pushout of (ADC4.1) along \(Q'\to Q'/S_r^Q\):
send a pushout class \([(b,[\lambda])]\) to \([b+j\lambda]\), and send
\([b]\) back to \([(b,0)]\). Changing \(\lambda\) by an element of
\(S_r^Q\) or using the relation \((j\mu,-[\mu])\) changes the first
formula by an element of \(S_r\). Both compositions are identity. The
formulas descend from continuous maps and are continuous for quotient
topologies.

The sheaf comparison retaining this row is RGR9's actual coefficient
fibre \[
\mathscr K_r=\operatorname{Cone}(-\mathrm{Gys}_{Q'}\,\sigma_r)[-1],
\qquad
\mathscr K_r\longrightarrow\mathscr E_r
\longrightarrow\underline{J'}[1]\longrightarrow\mathscr K_r[1].
\tag{ADC4.3}
\] The first map uses \(j\) on nearby coefficients and \(1_C\) on
vanishing coefficients; the second uses \(k'\) and zero. On costalks
they are \[
[C\xrightarrow{\sigma_r}Q']\xrightarrow{(1_C,j)}
[C\xrightarrow{a_r}B]\longrightarrow J'[-1].
\tag{ADC4.4}
\] These are literal continuous current maps induced by the original
\(q\) and \(k\). Their cokernel row is (ADC4.2). In particular the
source of the arithmetic boundary-image part is the dual of the same
original LNC quotient \(Q\), rather than an unrelated coefficient.

The reduced arithmetic receiver \(\mathscr K_r\) has the endpoint-free
part of ADC3's cross, with \(B\) replaced by \(Q'\), \(a_r\) by
\(\sigma_r\), and boundary image \(Q'/S_r^Q\). Alternatively
\(\mathscr K_{r,\mathrm{full}}=\mathscr K_r\oplus i_*E_p'[-1]\) retains
exactly the endpoint terms in the full ADC3 table. Its
coinvariant-boundary kernel is exactly \(S_r^Q\). Both invariant-cycle
quotients still vanish. This isolates the original arithmetic
coefficient inside the full \(A'\) receiver without changing the test
for the attaching map.

Finally, the actual morphism \(\mathscr E_r\to D_cP\) of RGR5 has nearby
map \(1_B\) and vanishing map \(\sigma_r:C\to Q'\); on costalks it is \[
[C\xrightarrow{a_r}B]\xrightarrow{(\sigma_r,1_B)}
[Q'\xrightarrow jB].
\tag{ADC4.5}
\] Its cone is \(i_*\operatorname{Cone}(\sigma_r)\), retaining the
kernel \(C_{\rm crit}\) and quotient \(Q'/S_r^Q\). It connects ADC3 to
the original LNC/LVD dual local arrows by an actual sheaf map. No
equality of a Gysin Hom class and an \(M\)-module Ext class is inferred.

\subsection{ADC5. The coefficient covariance and the full image
map}\label{adc5.-the-coefficient-covariance-and-the-full-image-map}

Write \(T_t^Q\) for the actual dilation induced on \(Q\) and
\(U_t^Q=tT_{1/t}^Q\). Their value maps intertwine with the corresponding
bounded \(T_t,U_t\) on \(H\). GTAH's exact identities are \[
T_t^*=\mathsf JU_t\mathsf J,\qquad
U_t^*=\mathsf JT_t\mathsf J,\qquad t>0.
\tag{ADC5.1}
\] Direct evaluation gives, for \(F\in Q\), \(y\in H\), \[
(T_t^Q)'A_H(y)(F)
=\langle T_t EF,\mathsf Jy\rangle_+
=\langle EF,T_t^*\mathsf Jy\rangle_+
=A_H(U_ty)(F).
\] The second identity follows by exchanging \(T_t,U_t\). Since \(D_r\)
is diagonal, it commutes with both. Thus, with linear maps on the
conjugate space, \[
\sigma_r\,\overline U_t=(T_t^Q)'\sigma_r,\qquad
\sigma_r\,\overline T_t=(U_t^Q)'\sigma_r,
\] \[
a_r\,\overline U_t=T_t'a_r,\qquad
a_r\,\overline T_t=U_t'a_r,\quad U_t'=tT_{1/t}'.
\tag{ADC5.2}
\] All maps are continuous; transposes are strong-continuous because
continuous primal maps carry bounded sets to bounded sets. Consequently
\(S_r^Q,S_r\), their exact quotients in ADC4, and \(\ker a_r\) are
stable under these actions. Their induced quotient actions are
continuous even when an image quotient is not separated.

This gives an exact description on each value coordinate. Let
\(e_\rho\in H\) be the vector of value one at \(\rho\) and zero
elsewhere, and let \(\delta_{\rho^\#,0}\in Q'\) be evaluation at
\(\rho^\#\). Then \[
\sigma_r(\overline e_\rho)
=m_\rho\,\overline{d_r(\rho)}\,\delta_{\rho^\#,0},\qquad
a_r(\overline e_\rho)
=m_\rho\,\overline{d_r(\rho)}\,q'\delta_{\rho^\#,0}.
\tag{ADC5.3}
\] The calculation is substitution into (ADC1.8) and
\(m_{\rho^\#}=m_\rho\). In particular, every actual off-critical value
functional lies in \(S_r^Q\); it vanishes from this formula exactly when
\(\rho\) is on the critical line. Finite-coordinate vectors are dense in
\(H\), and the image map is continuous, but this does not assert density
of finite primal full-jet blocks in \(Q\) or closedness of \(S_r\).

On a nonzero image line the two actions have exactly matching source and
target eigenvalues: \[
\begin{array}{c|c|c}
\text{action}&C/C_{\rm crit}\text{ at }\overline e_\rho
&S_r\text{ at }q'\delta_{\rho^\#,0}\\ \hline
\overline U_t\text{ / }T_t'&t^{1-\overline\rho}&t^{\rho^\#}\\
\overline T_t\text{ / }U_t'&t^{\overline\rho}&t^{1-\rho^\#}.
\end{array}
\tag{ADC5.4}
\] The entries are equal because \(\rho^\#=1-\overline\rho\). This is an
intertwining map on the actual entire source, with displayed formulas on
its value coordinates; it is not a proposed replacement arithmetic.

The continuous pairing with the original full quotient is equally exact.
Define \[
Z_r=\{F\in Q:\sigma_r(x)(F)=0\text{ for every }x\in C\}.
\tag{ADC5.5}
\] Using (ADC1.8) and the positive inner product gives \[
\sigma_r(\overline y)(F)
=\langle EF,\mathsf JD_ry\rangle_+
=\langle D_r^*\mathsf JEF,y\rangle_+,
\] \[
Z_r=\ker(D_r^*\mathsf JE)
=\{F\in Q:F(\rho)=0\text{ at every actual off-critical }\rho\}.
\tag{ADC5.6}
\] For the last equality, the scalar multiplying the reflected value is
nonzero exactly at an off-critical zero; reflection preserves that
subset. All higher jets remain in \(Z_r\) when the values vanish. The
pairing descends to \[
(Q/Z_r)\times(C/C_{\rm crit})\longrightarrow\mathbb C,\quad
([F],[\overline y])\longmapsto\sigma_r(\overline y)(F).
\tag{ADC5.7}
\] It is separately induced by the continuous original pairing. Its left
and right radicals are zero by (ADC5.6) and the injectivity calculation
(ADC1.10). This proves a concrete relation to the original \(Q\),
without asserting that either infinite-dimensional quotient is the
entire continuous dual of the other. Topological completions and image
closures are additional calculated receivers, not used to change
(ADC5.7).

\subsection{ADC6. What the genuine power-cover maps do to this
cross}\label{adc6.-what-the-genuine-power-cover-maps-do-to-this-cross}

For a recovered integer \(n\geq1\), the proper-current map over
\(b_n(z_p)=z_p^n\) is a map \(b_{n*}\mathscr E_r\to\mathscr E_r\). Its
literal nearby source is \(B^n[1]\), with cyclic permutation monodromy,
and its nearby map is \[
(b_0,\ldots,b_{n-1})\longmapsto T_n'\sum_{j=0}^{n-1}b_j.
\tag{ADC6.1}
\] To derive it, pull back the target universal cover: it has \(n\)
contractible inverse components. On each component the holomorphic local
isomorphism preserves the positive orientation and pushes its constant
current to \(T_n'b_j\); add the branches. This is the same full branch
calculation as LNC6.14 and retains all \(n\) components.

On nearby invariants the map is \(nT_n'\); on the angular coinvariants,
identified by sum, it is \(T_n'\). A proper map has a single
inverse-image pole here, so its actual stalk maps and positive costalk
maps are \[
\begin{array}{c|cc}
&\text{first displayed degree}&\text{second displayed degree}\\ \hline
i^*\mathscr E_r=B[1]\oplus C&nT_n'&\overline U_n\\
i^!\mathscr E_r=[C\to B]&\overline U_n&T_n'.
\end{array}
\tag{ADC6.2}
\] For the point-supported coordinate this uses (ADC5.2):
\(a_r\overline U_n=T_n'a_r\). For the current differential, pushforward
commutes with \(d\), and Dirac pushforward gives
\(b_{n*}\delta_p\lambda=\delta_pT_n'\lambda\). The angular current
pushes forward with degree one; its \(n\) local inverse branches each
contribute \(1/n\). This proves the entire chain map, not only its
induced eigenvalues.

The independently constructed weighted geometric trace on the original
primal sheaf has transpose giving the reverse map
\(\mathscr E_r\to b_{n*}\mathscr E_r\). Its full nearby map and
invariant/coinvariant factors are \[
b\longmapsto(T_{1/n}'b,\ldots,T_{1/n}'b),\qquad
B_{\rm inv}:T_{1/n}',\qquad B_{\rm coinv}:nT_{1/n}'=U_n'.
\tag{ADC6.3}
\] The finite-coordinate diagonal has sum \(nT_{1/n}'\), which fixes the
angular factor. The point map is \(\overline T_n\), and the costalk
identity is exactly \(a_r\overline T_n=U_n'a_r\). Thus on the cross the
resulting maps are \[
i^*: (T_{1/n}',\overline T_n),\qquad
i^!: (\overline T_n,nT_{1/n}').
\tag{ADC6.4}
\] The existence of the entire sheaf trace, including its
finite-coordinate nearby/vanishing correction, is constructed in
DCA3--DCA7 of \nolinkurl{CC_ADJOINT_DEFECT_COVER_ACTIONS.md.} Its actual
remaining coordinates can be displayed here. Put \(V_n=B^n\), let
\(\mathsf t(v)_j=v_{j-1}\) with cyclic indices, and let
\(\Delta b=(b,\ldots,b)\), \(\Sigma_n v=\sum_jv_j\),
\(e_0b=(b,0,\ldots,0)\). The letter \(\mathsf t\) denotes sheet
monodromy and is distinct from the support \(\tau\). Then \[
\Phi(b_{n*}\mathscr E_r)=C\oplus(V_n/\Delta B),\quad
\mathrm{can}_n(v)=(0,[v]),\quad
\mathrm{var}_n(x,[v])=e_0a_rx-(\mathsf t-1)v,
\] \[
\mathsf t_\Phi(x,[v])=(x,[\mathsf t v-e_0a_rx]).
\tag{ADC6.6}
\] These are DCA3.7--DCA3.9's actual cut-cone coordinates. In particular
\(\mathrm{var}_n\mathrm{can}_n=1-\mathsf t\) and
\(\mathrm{can}_n\mathrm{var}_n=1-\mathsf t_\Phi\). The constant stalk
restriction is \(\Delta\), and its angular restriction is \(e_0a_r\);
applying \(\Sigma_n\) recovers exactly \(a_r\), with no lost degree
factor.

Proper-current pushforward on this full vanishing space is
\((x,[v])\mapsto\overline U_nx\). For the reverse map define \[
c_x=T_{1/n}'a_rx,\qquad
w_n(x)=(0,c_x,2c_x,\ldots,(n-1)c_x).
\] Its full vanishing arrow is \[
x\longmapsto(\overline T_nx,[w_n(x)]),\qquad
(\mathsf t-1)w_n(x)=ne_0c_x-\Delta c_x.
\tag{ADC6.7}
\] The zeroth component of the last equality is \((n-1)c_x\), and each
remaining component is \(-c_x\). Since \(a_r\overline T_nx=nc_x\), this
proves the entire variation square: \[
\mathrm{var}_n(\overline T_nx,[w_n(x)])
=\Delta T_{1/n}'a_rx.
\tag{ADC6.8}
\] The correction is also the actual gluing homotopy. In the punctured
circle complex \([V_n\xrightarrow{\mathsf t-1}V_n]\) in degrees
\((-1,0)\), the difference between the two stalk-restriction composites
is \(d\,w_n+w_n\,d\). DCA3.5 realizes the sheaf from the homotopy fibre
of \(Rj_*L\oplus i_*S^{\mathrm{st}}\to i_*i^*Rj_*L\),
\((u,s)\mapsto u-\mathrm{res}(s)\). On its restriction fibre the map is
explicitly \((s,v)\mapsto(f_Ss,f_Nv+w_n(s))\); substitution in
\(d(s,v)=(ds,\mathrm{res}(s)-dv)\) proves the chain identity. This
constructs the sheaf map rather than inferring it from cohomology. All
finite-coordinate operations are continuous. Thus (ADC6.3)--(ADC6.4) are
the invariant/coinvariant maps of this full arrow, not a holomorphic
\(n\)-th-root inverse. The additional \(V_n/\Delta B\) has been
retained.

Both maps descend to \(N_r=B/S_r\), by (ADC5.2). On the positive
boundary \(\partial_0:B\to N_r\oplus E_p'\) they commute because the
induced quotient actions are precisely the actions of \(T_n'\) or
\(nT_{1/n}'\) on the same representative \(b\). The coefficient
covariance also proves commutation with every map in (ADC4.2)--(ADC4.5).
The endpoint summands are transported by their original dual pole
characters: current pushforward uses the dual endpoint part of
\(\rho_p(n)\), and the weighted-transfer transpose uses the
corresponding \(n\rho_p(1/n)'\). Their boundary component remains zero
and their inclusion/projection is unchanged.

In the dual coordinate labels of (ADC2.10), these endpoint maps are
explicitly \[
\begin{array}{c|cc}
&E_+'&E_-'\\ \hline
\text{current pushforward}&\operatorname{diag}(1,n)&\operatorname{diag}(n,1)\\
\text{weighted-transfer transpose}&\operatorname{diag}(n,1)&\operatorname{diag}(1,n).
\end{array}
\tag{ADC6.5}
\] The second row is obtained by substituting \(t=1/n\) in (ADC2.10) and
multiplying the transpose by the retained degree \(n\).

For \(t>0\) not an integer, (ADC5.2) still gives the actual coefficient
actions. The maps in this paragraph with geometric degree \(n\) are not
reinterpreted as noninteger covers.

\subsection{ADC7. Testing Deligne's two weight directions on the actual
boundary}\label{adc7.-testing-delignes-two-weight-directions-on-the-actual-boundary}

DC6--DC8 use two independent inputs: an upper bound on the special-fibre
classes and a lower bound on the coinvariants and supported cohomology,
the latter obtained by genuine duality with its retained Tate factor.
The full source conclusion is stronger than (ADC3.4): it identifies the
low-weight piece with \(\ker\partial\), disjoint from the coinvariant
summand, and proves \(\partial j:K_i\to\operatorname{im}\partial_i\)
injective.

Here all these maps have been computed. In degree zero, special classes
map by \(a_r:C\to K_0=B\). Thus a separation that killed the
intersection of special image and coinvariants would kill exactly
\(S_r\), by (ADC3.6). We now test the existing actions on that
intersection.

For a fixed geometric degree \(n>1\), define the logarithmic weight of a
scalar \(\lambda\ne0\) as \(w_n(\lambda)=2\log|\lambda|/\log n\). This
is an explicit modulus calculation, not an assertion that the
coefficient space is a finite-dimensional mixed Frobenius module. Under
current pushforward, the invariant nearby line indexed by \(\rho^\#\)
has scalar \(n^{1+\rho^\#}\), whereas its angular coinvariant has scalar
\(n^{\rho^\#}\). Their respective logarithmic weights are \[
2(2-\Re\rho),\qquad 2(1-\Re\rho).
\tag{ADC7.1}
\] The decrease by two is the trace direction of this current
pushforward; it is not Deligne's pullback Tate increase.

For the weighted-transfer transpose, the invariant line has scalar
\(n^{-\rho^\#}\) and the angular line has scalar \(n^{1-\rho^\#}\).
Their weights are \[
-2(1-\Re\rho),\qquad 2\Re\rho.
\tag{ADC7.2}
\] This has the positive angular increase by two. Nevertheless the
special source class \(\overline e_\rho\) has exactly the same angular
scalar \(n^{\overline\rho}=n^{1-\rho^\#}\), by (ADC5.4). Under the
pushforward it likewise has exactly the same scalar as the angular
target, \(n^{\rho^\#}\). Consequently neither actual direction separates
the source quotient \(C/C_{\rm crit}\) from \(S_r\): on every nonzero
image line they have identical eigenvalues. Multiplying both sides of an
equivariant map by a common Tate character shifts both weights equally
and does not change this equality. Multiplying only one side would
require a different map or action; none is inserted here.

The full boundary image \(N_r\) has a further exact source component
that also prevents assigning it a uniform one-sided weight bound from
these formulas. For any \(s_0\) not an original nontrivial zero,
evaluation \[
\epsilon_{s_0}(a)=\Theta a(s_0)\in B
\tag{ADC7.3}
\] is continuous and has nonzero restriction to \(J\), because
\(F_0(s_0)\ne0\). Since \(S_r\subset J^\perp\), its class in \(N_r\) is
nonzero. Under current pushforward angular action \(T_n'\) it has scalar
\(n^{s_0}\); under weighted-transfer angular action \(nT_{1/n}'\) it has
scalar \(n^{1-s_0}\). The resulting weights \[
2\Re s_0,\qquad 2(1-\Re s_0)
\tag{ADC7.4}
\] are unbounded in both directions as \(s_0\) ranges outside the
original divisor. This uses the actual \(J'\) quotient in (ADC4.2), with
the full entire \(F_0\), including its nonzero endpoint and trivial-zero
values. It is not a hypothetical zeta spectrum.

Passing to the arithmetic subreceiver \(\mathscr K_r\) removes precisely
this \(J'\) quotient by the proven triangle (ADC4.3). It leaves the
actual row \[
0\to S_r^Q\to Q'\to Q'/S_r^Q\to0
\tag{ADC7.5}
\] and the special-image map \(\sigma_r\). Its kernel is still
\(C_{\rm crit}\), and its nonzero image lines still have the matching
weights (ADC5.4). Thus the next useful attempted reduction has been
performed: the full unwanted coefficient quotient is removed by an exact
existing map, while the arithmetic intersection that DC's stronger
argument would have to kill remains explicitly \(S_r^Q\).

Finally, the normal separator \(c\) from GMS/FOD has identity action on
the entire original \(\mathcal Q\); its transpose is identity on \(Q'\).
Hence it acts as identity on \(S_r^Q\), not zero. The actual equality
\(ce_+=0\) for the translated normal extension cannot therefore be
reused as \(c\sigma_r=0\). The precise relation is the original dual
extension and its pushout in ADC4. This verifies the existing attempt
using the full earlier results rather than substituting the already
killed normal class for the present attaching map.

\subsection{ADC8. The continuous transpose and the actual
original-quotient
cross}\label{adc8.-the-continuous-transpose-and-the-actual-original-quotient-cross}

The coefficient transpose can be computed without assuming a duality
theorem. Identify \[
C'=(\overline H)'\simeq H,\qquad
x\longmapsto(\overline y\longmapsto\langle x,y\rangle_+).
\tag{ADC8.1}
\] This is the strong topological Riesz isomorphism for the conjugate
Hilbert space. Evaluation \(A\to(A'_\beta)'\), followed by \(a_r'\),
gives \[
b_r:A\longrightarrow H,\qquad
b_r=D_r^*\mathsf J E q.
\tag{ADC8.2}
\] Indeed, on \(\overline y\), \[
a_r'(\mathrm{ev}_f)(\overline y)
=a_r(\overline y)(f)
=\langle E qf,\mathsf JD_ry\rangle_+
=\langle D_r^*\mathsf JE qf,y\rangle_+.
\tag{ADC8.3}
\] Every map here is continuous, by (ADC1.9). Define also \[
\beta_r:Q\to H,\qquad \beta_r=D_r^*\mathsf JE,\qquad
b_r=\beta_rq.
\tag{ADC8.4}
\] Its kernel is exactly \(Z_r\) from (ADC5.5)--(ADC5.6); it retains all
higher jets and all critical-line values.

The actual coefficient cone to which these maps apply is \[
\mathscr P_r=\operatorname{Cone}
(\underline A\xrightarrow{b_r\,\mathrm{ev}_p}i_*H)
=[\,\underline A^{-1}\xrightarrow{+b_r}i_*H^0\,].
\tag{ADC8.5}
\] This is constructed by the displayed map on the original source. The
following cross uses its literal two-term complex, so no arbitrary
exactness of continuous dualization is being presumed.

The required current-level comparison is now explicit. ACD2--ACD4 in
\nolinkurl{CC_ADJOINT_DEFECT_CONTINUOUS_DUAL.md} construct the fine
resolution \[
\mathscr L_r^{-1}=\mathscr E_A^0,\quad
\mathscr L_r^0=\mathscr E_A^1\oplus i_*H,\quad
\mathscr L_r^1=\mathscr E_A^2,\qquad
d^{-1}f=(-df,b_rf(p)),\quad d^0(\omega,h)=-d\omega.
\tag{ADC8.5a}
\] Inclusion of locally constant functions and the identity on \(H\)
compare (ADC8.5) with this resolution. The coefficient-valued Poincaré
contraction proves this is a quasi-isomorphism. It does not require the
image of \(b_r\) to be closed. Use the fixed-support smooth seminorms
and their compact-support locally convex inductive limit to form the
actual continuous Hom sheaf
\(\mathscr D_r^k(U)=\Gamma_c(U,\mathscr L_r^{-k})'\), with
\(d_{\mathscr D}\ell=(-1)^{k+1}\ell d_{\mathscr L}\). Its degrees
\((-1,0,1)\) are the current terms
\(\mathscr T_A^{-2},\mathscr T_A^{-1}\oplus i_*H',\mathscr T_A^0\), and
\[
d_{\mathscr D}^{-1}v=(d_Tv,0),\qquad
d_{\mathscr D}^{0}(u,\lambda)=d_Tu-\delta_pb_r'\lambda.
\tag{ADC8.5b}
\] For instance the second formula on a test \(f\) is
\(-u(-df)-\lambda(b_rf(p))\); this fixes both signs. With \(R:C\to H'\),
\(R(\overline y)(h)=\langle h,y\rangle_+\), equation (ADC8.3) gives
\(b_r'R=a_r\). The map \[
\mathscr D_r\longrightarrow\mathscr E_r,\qquad
v\mapsto-v\ (\deg-1),\quad
(u,Rx)\mapsto(u,x)\ (\deg0),\quad
w\mapsto-w\ (\deg1)
\tag{ADC8.5c}
\] is a continuous chain isomorphism: the differential composites are
\((d_Tv,0)\) and \(-d_Tu+\delta_pa_rx\). This identifies
\(\mathscr E_r\) with the literal continuous support-dual complex of the
constructed \(\mathscr P_r\) resolution. It is not an assertion that
continuous dualization is exact on every locally convex complex. The
positive compact-form pairing in outer degrees has the minus signs just
displayed; the point pairing is \(R(x)(h)\). ACD3's compact-support
contraction proves the local support comparison, with the same positive
Dirac class as (ADC2.8).

Its restriction to the puncture is \(\underline A[1]\), and its pole
complex is \([A^{-1}\xrightarrow{b_r}H^0]\). Therefore \[
R\psi\mathscr P_r=A[1],\quad T=1,\quad
H^{-1}(i^*\mathscr P_r)=\ker b_r,\quad
H^0(i^*\mathscr P_r)=H/b_rA.
\tag{ADC8.6}
\] The map to \(A[1]\oplus A[\vartheta_p]\) is identity on the
degree-minus-one \(A\) and zero on \(H\). In its restriction fibre the
degree-zero coordinate is \(H\oplus A\), with preceding differential
\(a\mapsto(b_ra,a)\). Its quotient coordinate is
\((h,a)\mapsto h-b_ra\). This cancels the identity pair and gives \[
i^!\mathscr P_r\simeq H[0]\oplus A[-1],\qquad
\partial_{-1}=-b_r:A\to H,\qquad
\partial_0=+1:A[\vartheta_p]\to A.
\tag{ADC8.7}
\] The last sign is positive after the same cochain shift as in
LVD3/RGR6; it also follows from the retained positive Dirac coordinate
in the fibre. The full source includes the endpoint term \[
\mathscr P_{r,\mathrm{full}}=\mathscr P_r\oplus i_*E_p''[1],
\tag{ADC8.8}
\] where \(E_p''\simeq E_p\) by finite-dimensional evaluation, with both
original labels retained. This contributes special/support degree
\(-1\), no nearby group, and unshifted vanishing degree \(-2\).

The resulting entire nonzero cross is:

{\def\LTcaptype{none} % do not increment counter
\begin{longtable}[]{@{}
  >{\raggedright\arraybackslash}p{(\linewidth - 14\tabcolsep) * \real{0.1250}}
  >{\raggedright\arraybackslash}p{(\linewidth - 14\tabcolsep) * \real{0.1250}}
  >{\raggedright\arraybackslash}p{(\linewidth - 14\tabcolsep) * \real{0.1250}}
  >{\raggedright\arraybackslash}p{(\linewidth - 14\tabcolsep) * \real{0.1250}}
  >{\raggedright\arraybackslash}p{(\linewidth - 14\tabcolsep) * \real{0.1250}}
  >{\raggedright\arraybackslash}p{(\linewidth - 14\tabcolsep) * \real{0.1250}}
  >{\raggedright\arraybackslash}p{(\linewidth - 14\tabcolsep) * \real{0.1250}}
  >{\raggedright\arraybackslash}p{(\linewidth - 14\tabcolsep) * \real{0.1250}}@{}}
\toprule\noalign{}
\begin{minipage}[b]{\linewidth}\raggedright
\(i\)
\end{minipage} & \begin{minipage}[b]{\linewidth}\raggedright
special \(A_i\)
\end{minipage} & \begin{minipage}[b]{\linewidth}\raggedright
\(K_i\)
\end{minipage} & \begin{minipage}[b]{\linewidth}\raggedright
punctured \(B_i\)
\end{minipage} & \begin{minipage}[b]{\linewidth}\raggedright
nearby invariants \(C_i\)
\end{minipage} & \begin{minipage}[b]{\linewidth}\raggedright
support \(O_i\)
\end{minipage} & \begin{minipage}[b]{\linewidth}\raggedright
\(\alpha_i\)
\end{minipage} & \begin{minipage}[b]{\linewidth}\raggedright
\(\partial_i\)
\end{minipage} \\
\midrule\noalign{}
\endhead
\bottomrule\noalign{}
\endlastfoot
\(-2\) & \(0\) & \(0\) & \(0\) & \(0\) & \(E_p''\) & \(0\) & \(0\) \\
\(-1\) & \(\ker b_r\oplus E_p''\) & \(0\) & \(A\) & \(A\) & \(H\) &
\((a,e)\mapsto a\) & \(-b_r\) \\
\(0\) & \(H/b_rA\) & \(A[\vartheta_p]\) & \(A[\vartheta_p]\) & \(0\) &
\(A\) & \(0\) & \(+1\) \\
\end{longtable}
}

In particular the genuine invariant-cycle obstruction is now \[
\boxed{\operatorname{im}\partial_{-1}/\partial_{-1}j_{-1}K_{-1}
=b_rA=\beta_rQ,\qquad
\operatorname{coker}\operatorname{sp}_{-1}
=A/\ker b_r\simeq Q/Z_r.}
\tag{ADC8.9}
\] The first identification of the cokernel with the image uses
\([a]\mapsto-b_ra\), retaining the boundary sign. The induced algebraic
bijection \(Q/Z_r\to b_rA\) can equivalently use \([F]\mapsto\beta_rF\)
when the minus is stated separately. Quotient topologies are retained,
and no inverse for its ambient image topology is asserted. All other
obstruction images vanish.

This dual direction receives the original LNC quotient through an
explicit map of the full local source. Put
\(\mathsf P=\mathscr F_{\rm red}[1]\simeq\operatorname{Cone}(q\,\mathrm{ev}_p)\),
with terms \(\underline A^{-1}\to i_*Q^0\). Then \[
f_r:\mathsf P\longrightarrow\mathscr P_r,\qquad
f_r^{-1}=1_A,\quad f_r^0=\beta_r
\tag{ADC8.10}
\] is a continuous sheaf chain map since \(\beta_rq=b_r\). On the full
objects add the original evaluation \(E_p[1]\to E_p''[1]\). Its nearby
map is identity \(A\), and its vanishing map is \(\beta_r:Q\to H\). On
support degree zero it is \(\beta_r\), and on support degree one it is
identity \(A\). Thus the original boundary \(-q:A\to Q\) and the new
boundary \(-b_r:A\to H\) satisfy \[
\begin{array}{ccc}
A&\xrightarrow{-q}&Q\\
\Vert&&\downarrow\beta_r\\
A&\xrightarrow{-b_r}&H .
\end{array}
\tag{ADC8.11}
\] It satisfies all four commuting-map identities (DC10.2), with the
explicit shifts and coefficients just computed. Here those identities
are proved on the displayed continuous complexes; the finite-dimensional
Frobenius hypotheses used for DC10's weight conclusions are not asserted
for them. Its induced map on obstruction images is the actual surjection
\[
Q\xrightarrow{\beta_r}\beta_rQ,\qquad
\ker\beta_r=Z_r.
\tag{ADC8.12}
\] This constructs the precise arithmetic subquotient of the original
local defect, rather than naming two different receivers as if they were
identical.

The transpose is also the actual old current map (ADC4.5), with its
representative sign retained. On the fine resolutions of the plus-\(q\)
and plus-\(b_r\) cones, \(f_r\) is identity on forms and \(\beta_r\) on
the point term; its transpose is identity on currents and
\(\beta_r'R=\sigma_r\) on point functionals. The literal dual of the
plus-\(q\) cone in the sign convention (ADC8.5c) is
\(\operatorname{Cone}(-\delta_pq')[-1]\). The accepted RGR
representative is \(\operatorname{Cone}(+\delta_pq')[-1]\); their
comparison negates the point coordinate. Therefore the map into that
accepted representative is \((u,x)\mapsto(u,-\sigma_rx)\), as proved
directly in (ACD4.4). Its positive costalk coordinates are still
\((\sigma_r,1_B)\), by the simultaneous coordinate comparisons in
(RGR6.3a). Thus the local source square and the earlier current
comparison are related by an actual continuous transpose, rather than by
a coincidence of coefficient groups.

Finally \(b_r=0\) if and only if \(D_r=0\): if \(b_r=0\), the bounded
map \(D_r^*\mathsf J\) vanishes on dense \(E(Q)\), hence on \(H\);
conversely \(D_r=0\) makes every formula zero. Therefore the new cross's
surjectivity in degree \(-1\) is equivalent to vanishing of the actual
adjoint discrepancy. Unlike (ADC3.4), its obstruction image is the part
that detects that discrepancy.

The exact coefficient actions on this new boundary can also be
calculated: \[
\beta_r T_t^Q=U_t^*\beta_r,\qquad
\beta_r U_t^Q=T_t^*\beta_r,\qquad
b_rT_t^A=U_t^*b_r,\quad b_rU_t^A=T_t^*b_r.
\tag{ADC8.13}
\] For the first equality, \(E T_t^Q=T_tE\) and
\(\mathsf JT_t=U_t^*\mathsf J\) by (ADC5.1); the diagonal \(D_r^*\)
commutes with \(U_t^*\). This gives
\(D_r^*\mathsf J T_tE=U_t^*D_r^*\mathsf JE\). The second follows with
\(T_t,U_t\) exchanged; composing with \(q\) gives the remaining
formulas. They define actual continuous chain maps on (ADC8.5), and show
that the two coefficient directions descend to its invariant-cycle
image. For a value class in \(Q/Z_r\) supported at an actual \(\rho\),
its nonzero image lies at \(\rho^\#\) in \(H\). A representative in the
complete \(Q\) can still have higher-jet Jordan terms; these lie in the
recorded kernel \(Z_r\), rather than being deleted from \(Q\). The
source and target scalars on the indicated quotient and image are
respectively \[
t^\rho=t^{\,1-\overline{\rho^\#}},\qquad
t^{1-\rho}=t^{\,\overline{\rho^\#}}.
\tag{ADC8.14}
\] Their weights again agree on every nonzero image line. This proves
the compatibility directly on the new cross, rather than inferring a
vanished map from the existence of the two actions.

\subsection{ADC9. Exact outcome and source-use
record}\label{adc9.-exact-outcome-and-source-use-record}

The two calculated crosses test different parts of the same actual
attaching data: \[
\begin{array}{c|c|c}
\text{receiver}&\text{invariant-cycle obstruction image}
&\text{additional retained defect}\\ \hline
\mathscr E_r&0&S_r=\ker(K_0\to\operatorname{im}\partial_0)\\
\mathscr P_r&b_rA\simeq Q/Z_r&
H/b_rA\text{ in special degree }0.
\end{array}
\tag{ADC9.1}
\] Their source maps are the actual \(a_r(\overline y)=q'A_H(D_ry)\) and
\(b_r=D_r^*\mathsf JE q\), not independent examples. The complete
point-source quotient \(C/C_{\rm crit}\) and the original arithmetic
quotient \(Q/Z_r\) are joined by the nondegenerate pairing (ADC5.7). The
comparison maps with the original source are (ADC4.5), (ADC8.10), and
the boundary square (ADC8.11). The full original dual extension survives
as the pushout (ADC4.2), with the source kernel and every endpoint
degree retained.

The next attempted step required by WU065 has therefore been carried
out. Starting from the automatically vanishing invariant-cycle cokernel
of \(\mathscr E_r\), the calculation first isolates the stronger
coinvariant-injectivity defect, then removes its full \(J'\) coefficient
quotient by the actual source triangle, and then constructs the
transpose cone receiving the original invariant-cycle quotient. Its
obstruction is precisely the image of the adjoint discrepancy on
original classes. The two actual coefficient directions are computed on
that image in (ADC8.13)--(ADC8.14); neither has separated source and
target weights there. This describes the proved effect of the existing
maps, without declaring a finite-field purity theorem for these
coefficients or using a new twist to force one.

Reading coverage for this derivation: the connected
\nolinkurl{USER_ARGUMENT_RECONSTRUCTION.md,} its correction-precedence
table, the verbatim WU062 and WU064--WU065 passages; the complete
DC0--DC12 proof; LVD0--LVD9; RGR0--RGR12, including the full
original-source, current-cone and positive-adjoint distinctions;
GTAH0--GTAH4; ACD0--ACD8, with the literal current-resolution and
source-map comparison in ACD2--ACD4 used in ADC8; DCA0--DCA12, including
its complete cut-cone, genuine cover/trace and composition homotopies,
with DCA3 and DCA6 used explicitly in ADC6; and the accepted LNC0--LNC11
proof, previously authored and independently reviewed here. LNC's
current, nearby and support maps are used with their already proved
continuous contractions and signs. The original source references and
their exact witness status are preserved in DC0, LNC11, LVD0 and RGR0.
In particular Deligne's local invariant-cycle proof is read through the
identified current French transcription of Weil II; this is not claimed
to be Deligne-authored TeX.

This is a calculation on the actual recovered arithmetic and
current/sheaf maps. It establishes neither an off-critical zero nor the
vanishing of \(D_r\). It supplies the exact subquotients and commuting
maps through which that discrepancy enters the stronger Deligne
comparison.
